//323111047

\section{Introducción}

Planteamiento del problema. El reto técnico de esta actividad consiste en realizar una transición hacia el lenguaje de programación Python para implementar diversas tareas de lógica de programación que abarcan desde el manejo de tipos de datos básicos hasta la manipulación de estructuras de datos complejas. El problema se centra en dominar la sintaxis de un lenguaje interpretado y dinámicamente tipado, resolviendo nueve ejercicios específicos que incluyen operaciones aritméticas, gestión de archivos, procesamiento de cadenas de texto y el uso de bibliotecas especializadas para funciones matemáticas. Además, se requiere estructurar la ejecución del programa mediante menús interactivos y funciones que permitan organizar el código de manera modular dentro de un entorno de celdas de código y texto.

\vspace{0.4cm}

Motivación. La importancia de realizar este conjunto de ejercicios radica en la relevancia de Python como herramienta fundamental en la ingeniería moderna, especialmente en áreas como el análisis de datos, la automatización y el desarrollo de aplicaciones científicas. Comprender la flexibilidad de sus estructuras de datos, como las pilas, colas y diccionarios, así como la potencia de sus mecanismos de entrada y salida, permite al estudiante desarrollar soluciones más legibles y eficientes en comparación con lenguajes de bajo nivel. Esta práctica sirve como base para abordar problemas de mayor escala que requieren una alta capacidad de abstracción y persistencia de datos mediante el manejo de archivos.

\vspace{0.4cm}

Objetivos. Se busca que el alumno demuestre dominio en el uso de operadores aritméticos, manipulación de archivos y aplicación de funciones matemáticas y de procesamiento de texto mediante f-strings y técnicas de slicing. Asimismo, se pretende validar la correcta organización de la información mediante estructuras de datos lineales y diccionarios, integrando funciones normales y anónimas dentro de una lógica de control de flujo, lo que permitirá demostrar la versatilidad del lenguaje en la resolución de problemas algorítmicos.